\documentclass[11pt,a4paper]{article}

\usepackage[T1]{fontenc}
\usepackage[utf8]{inputenc}
\usepackage[margin=22mm]{geometry}
\usepackage{booktabs}
\usepackage{array}
\usepackage{enumitem}
\usepackage{titlesec}
\usepackage{fancyhdr}
\usepackage{listings}
\usepackage[hidelinks]{hyperref}

\pagestyle{fancy}
\fancyhf{}
\lhead{\small BeamMP Race Manager}
\rhead{\small Phase 6}
\cfoot{\small \thepage}
\renewcommand{\headrulewidth}{0.4pt}

\titleformat{\section}{\large\bfseries}{\thesection}{0.6em}{}
\titlespacing{\section}{0pt}{14pt}{6pt}

\setlist[itemize]{leftmargin=*,itemsep=2pt,topsep=4pt}
\setlist[enumerate]{leftmargin=*,itemsep=2pt,topsep=4pt}
\setlength{\parskip}{6pt}
\setlength{\parindent}{0pt}

\lstset{
  basicstyle=\ttfamily\scriptsize,
  frame=single,
  framesep=5pt,
  columns=fullflexible,
  keepspaces=true,
  breaklines=true,
  showstringspaces=false,
  xleftmargin=0pt,
  aboveskip=8pt,
  belowskip=8pt,
}

\newcommand{\code}[1]{\texttt{\small #1}}

\begin{document}

\begin{center}
{\LARGE\bfseries Phase 6: what was built}\\[4pt]
{\small BeamMP Race Manager \quad Client: dhardesty99}
\end{center}

\vspace{4pt}
\hrule
\vspace{8pt}

Phase 6 is on the server. It is the last phase of the plan: CoPilot and
Chase, the challenges, the language switch and the tracking switch. This
is what each part does and how to use it. The build is 0.7.5.

\section{CoPilot and Chase}

Watch another driver from their car. There are two ways in, and both sides
have to agree.

\begin{itemize}
  \item \textbf{Invite.} You are driving and want somebody to ride with
        you. Open CoPilot, press Invite next to their name. They accept or
        say no.
  \item \textbf{Ask to watch.} You want to watch a driver. Press Ask to
        watch next to their name. They accept or say no.
\end{itemize}

On yes, the watcher's camera goes into the driver's car. The C key changes
the camera as usual, so CoPilot is the view from the seat and Chase is the
view from behind. The watcher's own car stays where it was parked, and
their own timing is not touched.

\begin{itemize}
  \item Stop watching brings you back to your own car. A driver can also
        send everybody watching them back with one button.
  \item The plain TAB switch is shut on the server. A switch onto anybody
        else's car is undone at once, so watching is by invite or request
        only.
  \item Arming a run ends watching, because your own car is the one that
        races. Somebody mid run cannot be made to watch, and cannot ask
        to.
  \item If the driver changes car, the watcher is moved to the new one. If
        the driver leaves, the watcher is sent home. An offer nobody
        answers goes away after a minute.
\end{itemize}

\section{Challenges}

Daily and weekly time attacks, posted by admins, entered by everybody.

\subsection*{Posting one}

Open Challenges and press Post a challenge. Give it a name and a few lines
of description, pick daily or weekly, the course, the laps if the course
is a circuit, the classes that may enter, and the ladder of times: each
rung is a time to beat and the XP it pays, typed as \code{1:32.43} or
\code{92.43}. Up to twenty rungs. Pick when it starts: now, in one hour,
in six, or tomorrow. A daily runs for a day from its start, a weekly for
seven days.

Three daily and five weekly challenges can be open at once. Ending one, or
one running out, makes room. Edit changes any of it live; a change of
course or laps clears its board, because those times were driven on the
old one. End it now and Delete are there too.

\subsection*{Entering one}

The board shows daily challenges in one table and weekly in another, each
a set size that scrolls. Under each name, in grey, the classes and the
course. Click a challenge and the window opens out with everything about
it: the description, the course and laps, the classes, when it ends, the
ladder, your best time, and the board of everyone who has finished it.
Pick your class if it asks for one, controller or wheel, and press Enter
and go to the grid. From there it is a normal run.

\subsection*{Scoring}

Your best time on each challenge is kept. Finishing under a rung pays
that rung's XP, once. A faster run later pays only the difference up to
the higher rung. A slower run changes nothing. A time outside the ladder
still goes on the board. The results row of the run says which rung you
reached and what it paid.

\begin{itemize}
  \item Teams cannot enter a challenge. Break up the team first.
  \item Guests count, the same as everybody.
  \item A challenge that is over stays on the board a while, greyed, so
        the times can still be read.
\end{itemize}

\section{Options}

\begin{itemize}
  \item \textbf{Language.} English or Spanish. The bars, the window titles
        and the CoPilot and Challenges windows change. The rest stays in
        English.
  \item \textbf{XP and challenge tracking.} On or off for yourself. Off
        means a run counts for nothing: no XP, no challenge time. It is
        for practice. Turn it back on before a race.
\end{itemize}

\section{Guests and accounts}

Everybody on the server counts, guest or not: a name picked once, a level,
records, XP and challenges. A guest is recognised on the next visit by
their connection, and the six character code from the first join brings
their name back if that ever changes.

Whether guests may join at all follows Bobby's file on the server,
\code{Resources/Server/PatreonAuth/allowed\_discord\_ids.json}. While its
\code{enabled} is false, guests join as normal. When Bobby switches it on,
a guest is turned away at the door with a message to log in with a BeamMP
account, and accounts still get in. The file is read every ten seconds,
so the switch works live. On the console, \code{rm guests} says the
current state. Note that anyone who joins as a guest, including yourself,
is turned away while it is on.

Somebody logged into a BeamMP account is named after their account on
first join, with no card to fill in, as long as that name is one the
server can take and nobody has it. Names already picked are never
replaced.

\section{Discord}

The Open in browser button goes through the BeamMP launcher, which opens
the link in your browser. The game keeps focus, so Alt-Tab to it. Copy
link is still there.

\section{The console}

\begin{lstlisting}
rm challenges             the challenges on the board
rm challenge end <id>     end a challenge now
rm challenge delete <id>  take a challenge off the board
rm watching               who is watching whom
rm guests                 whether guests may join, and where that is decided
\end{lstlisting}

\section{Testing}

Every part of phase 6 runs against the same test host as the earlier
phases, without the game. The watching flow, the challenge board, the
ladder, the limits, the schedule, the tracking switch, guests and
accounts, and the camera side in the car are all covered: 166 checks in
three new suites, and every earlier suite still passes. The interface is
audited for the mistakes that only show in the game, and each fix was
checked against a broken version to be sure the test would catch it.

\end{document}
